%!TEX TS-program = Arara
% arara: pdflatex: {shell: no}
% arara: pdflatex: {shell: no}

\documentclass[12pt,parskip=half]{scrartcl}

\author{Uwe Ziegenhagen}
\title{Mein erstes \LaTeX-Dokument}

\usepackage[ngerman]{babel} % Silbentrennung
\usepackage{blindtext} % Dummy-Text
\usepackage{microtype} % Mikrotypografie-Erweiterungen aktivieren

\usepackage{graphicx} % Dateitypen: PNG, PDF, JPG

\usepackage{showlabels}
\usepackage{todonotes}

\usepackage{booktabs}
\usepackage{multirow}

\usepackage{amssymb}
\usepackage{nicefrac}
\usepackage{siunitx}

\usepackage[standard,thmmarks]{ntheorem}


\usepackage{hyperref} % Links, als letztes Paket laden
\usepackage{hyperref} % Link Setup
\hypersetup{
    bookmarks=true,                     % show bookmarks bar
    unicode=false,                      % non - Latin characters in Acrobat’s bookmarks
    pdftoolbar=true,                        % show Acrobat’s toolbar
    pdfmenubar=true,                        % show Acrobat’s menu
    pdffitwindow=false,                 % window fit to page when opened
    pdfstartview={FitH},                    % fits the width of the page to the window
    pdftitle={Mein erstes Dokument},                        % title
    pdfauthor={Uwe Ziegenhagen},                 % author
    pdfsubject={Abschlussarbeit},                   % subject of the document
    pdfcreator={LaTeX},                   % creator of the document
    pdfproducer={pdflatex},             % producer of the document
    pdfkeywords={Uwe, Ziegenhagen},   % list of keywords
    pdfnewwindow=true,                  % links in new window
    colorlinks=true,                        % false: boxed links; true: colored links
    linkcolor=blue,                          % color of internal links
    filecolor=blue,                     % color of file links
    citecolor=blue,                     % color of file links
    urlcolor=blue                        % color of external links
}


\begin{document}
\maketitle

\tableofcontents % Inhaltsverzeichnis

\listoffigures

\listoftables

\listoftodos

\clearpage

\section{ewrwerwer}

\section{Einführung}
\subsection{Fragestellung} \todo{fsadgfdsg}

In Abschnitt \ref{sec:fazit} zeige ich warum, in Abbildung \ref{fig:katze} auf Seite \pageref{fig:katze} zeige ich wieso

Wieviel \% brauche ich für 100 \$? % Diverse Sonderzeichen brauchen den Backslash

Es gibt Empfehlungen zu LaTeX.
Am besten pro Satz eine Zeile benutzen.
Dann tut sich LaTeX leichter damit, den Ort eines Fehlers zu kennzeichnen.

\includegraphics[width=0.75\textwidth]{Bilder/Katze}


\begin{figure}
\centering
\includegraphics[width=0.75\textwidth]{Bilder/Katze}
\caption{Meine freche Katze}\label{fig:katze}
\end{figure}


Ein neuer Absatz erfordert eine leere Zeile davor

\blindtext[5]

\subsection{Literaturüberblick}
\subsubsection{Deutschsprachige Literatur}

\blindtext[5]

\subsubsection{Internationale Literatur}

\blindtext[5]

% Caption geht in den KOMA Klassen auch über den 
% captionof Befehl
\begin{center}
\includegraphics[width=0.75\textwidth]{Bilder/Katze}
\captionof{figure}{Meine superfreche Katze}\label{fig:katze423}
\end{center}



\paragraph{Hallo} \blindtext[5]

\begin{figure}
\centering
\includegraphics[width=0.75\textwidth]{Bilder/Katze1}
\caption{Meine kuschelige Katze}\label{fig:katze1}
\end{figure}



\paragraph{Welt} \blindtext[5]

\subparagraph{Hello} \blindtext[5]


\section{Untersuchung}

\blindtext[5] 



\section{Fazit}\label{sec:fazit}

\begin{figure}
\centering
\includegraphics[width=0.75\textwidth,angle=45]{Bilder/miau}
\caption{Meine wuschelige Katze}\label{fig:miau}
\end{figure}


\blindtext[5]

\section{Tabellen}

Hallo, ich bin ein Text mit \textbf{verschiedenartig} \textit{formatierten}    Wörtern, dies \textbf{\textit{kann man auch}} kombinieren, einige \textsl{Fonts haben auch} geneigte Schnitte

\clearpage

\begin{tabular}{|l|c|r|p{5cm}|} \hline
\textbf{Spalte 1} & \textbf{Spalte 2} & \textbf{Spalte  3} & \textbf{Spalte  4} \\ \hline \hline
Hallo & Welt & ich & bin ein kleiner Text \\ \hline 
Hallihallo & Welt und World & ich bin nicht & bin ein kleiner Text, mit vielen Wörtern \\ \hline
\end{tabular}\vspace*{1cm}

\begin{table}
\caption{Meine erste schöne Tabelle}\label{tab:erste}
\begin{tabular}{lcrp{5cm}} \toprule
\textbf{Spalte 1} & \textbf{Spalte 2} & \textbf{Spalte  3} & \textbf{Spalte  4} \\ \midrule
Hallo & Welt & ich & bin ein kleiner Text \\ 
Hallihallo & Welt und World & ich bin nicht & bin ein kleiner Text, mit vielen Wörtern \\ \bottomrule
\end{tabular}
\end{table}

\begin{table}
\caption{Meine erste schöne Tabelle}\label{tab:zweite}
\begin{tabular}{lcrp{5cm}} \toprule[1.5pt]
\textbf{Spalte 1} & \textbf{Spalte 2} & \textbf{Spalte  3} & \textbf{Spalte  4} \\ \cmidrule[2pt](rl){1-4}
Hallo & Welt & ich & bin ein kleiner Text \\ 
Hallihallo & Welt und World & ich bin nicht & bin ein kleiner Text, mit vielen Wörtern \\ \bottomrule[1.5pt]
\end{tabular}
\end{table}

\begin{tabular}{ccccccccc}
1 & 2 & 3 & 4 & 5 & 6 & 7 & 8 & 9 \\
1 & 2 & 3 & 4 & 5 & 6 & 7 & 8 & 9 \\
1 & 2 & 3 & 4 & 5 & 6 & 7 & 8 & 9 \\
\multicolumn{2}{c}{a} & 3 & 4 & 5 & 6 & 7 & 8 & 9 \\
1 & 2 & 3 & 4 & 5 & 6 & 7 & 8 & 9 \\
1 & 2 & 3 & 4 & 5 & 6 & 7 & 8 & 9 \\
1 & 2 & 3 & 4 & 5 & 6 & 7 & 8 & 9 \\
\multirow[c]{2}{*}{a} & 2 & 3 & 4 & 5 & 6 & 7 & 8 & 9 \\
  & 2 & 3 & 4 & 5 & 6 & 7 & 8 & 9 \\
1 & 2 & 3 & 4 & 5 & 6 & 7 & 8 & 9 \\
1 & 2 & 3 & 4 & 5 & 6 & 7 & 8 & 9 \\
\end{tabular}

Hallo Welt

Hallo \hspace*{1cm}       Welt

Hallo \hfill Welt \vfill Welt

\section{Mathematik-Satz mit \LaTeX}

Hallo, ich bin eine Formel $a^2+b^2=c^2$ im Fließtext in \TeX-Notation.

Hallo, ich bin eine abgesetzte Formel $$a^2+b^2=c^2$$ in \TeX-Notation.

Hallo, ich bin eine Formel \(a^2+b^2=c^2\) im Fließtext in \LaTeX-Notation.

Hallo, ich bin eine abgesetzte Formel \[a^2+b^2=c^2\] in \LaTeX-Notation.

Siehe Formel \ref{eq:f1}

\begin{equation}
a^2+b^2=c^{-2} \Rightarrow \int_{\infty}^{-1}  \sin x  \alpha \beta \gamma \label{eq:f1}
\end{equation}

\begin{equation}
x_{1, 2} = -\frac{p}{2} \pm \sqrt{ \left(\frac{p}{2}\right)^2 - q }
\end{equation}

\begin{equation}
\sum_{x=1}^{\infty} x^2  \rightarrow \prod_{x=1}^{\infty} x^2  \Rightarrow \prod_{x=1}^{\infty} x^2 \qquad \nicefrac{1}{2}\Box
\end{equation}

% https://www.dickimaw-books.com/latex/thesis/html/ntheorem.html


\begin{Definition}[Tautology]\label{def:tautology}
A \emph{tautology} is a proposition that is always true for any value of its variables.
\end{Definition}

\begin{Definition}[Contradiction]\label{def:contradiction}
A \emph{contradiction} is a proposition that is always false for any value of its variables.
\end{Definition}

\begin{Theorem}
If proposition $P$ is a tautology then $\sim P$ is a contradiction, and conversely.
\begin{Proof}
If $P$ is a tautology, then all elements of its truth table are true (by Definition~\ref{def:tautology}), so all elements of the truth table for $\sim P$ are false, therefore $\sim P$ is a contradiction (by Definition~\ref{def:contradiction}).
\end{Proof}
\end{Theorem}

\begin{Example}\label{ex:rain}
``It is raining or it is not raining'' is a tautology, but ``it is not raining and it is raining'' is a contradiction.
\end{Example}

\begin{Remark}
Example~\ref{ex:rain} used De Morgan's Law $\sim (p \vee q) \equiv \sim p \wedge \sim q$.
\end{Remark}

\begin{Proof}
Example~\ref{ex:rain} used De Morgan's Law $\sim (p \vee q) \equiv \sim p \wedge \sim q$. \qed
\end{Proof}

\section{Einheitensatz mit \texttt{siunitx}}

\num{1234567890,15645}

$100 m/s^2$ sieht doof aus

\si{\m/\second^2}

\SI{1231,484810}{\m/\second^2}

\ang{90}

\numrange{1}{100}

\begin{tabular}{lS}
1 & 12345345.46 \\
1 & 1345.7 \\
1 & 12345345.897987 \\
\end{tabular}


\end{document}

